% =========================================================================
\section{Funciones de la posición relativa, delta de Dirac y desarrollo multipolar}\label{sec:posicion}
% =========================================================================
En electromagnetismo los campos suelen producirlos \emph{fuentes} (cargas,
corrientes) situadas en puntos $\vc{r}\,'$, y se evalúan en otros puntos
$\vc{r}$ donde queremos conocer el campo. Toda la dependencia aparece a través
del vector que une ambos puntos~\cite{FORMULARIO,GRIFFITHS}.

% -------------------------------------------------------------------------
\subsection{Punto fuente, punto campo y vector relativo}\label{sec:relativa}
% -------------------------------------------------------------------------
\begin{definicion}[Posición relativa]\label{def:relativa}
El \emph{punto fuente} es $\vc{r}\,'=(x',y',z')$ y el \emph{punto campo} es
$\vc{r}=(x,y,z)$. El vector que va de la fuente al punto campo, su módulo y su
versor son
\begin{align}
  \vc{R}&=\vc{r}-\vc{r}\,'=(x-x')\,\ux+(y-y')\,\uy+(z-z')\,\uz=R_x\,\ux+R_y\,\uy+R_z\,\uz,\label{eq:R-vec}\\
  R&=\norm{\vc{R}}=\bigl[(x-x')^{2}+(y-y')^{2}+(z-z')^{2}\bigr]^{1/2},\qquad
  \uR=\frac{\vc{R}}{R}\quad\Bigl(\vc{R}=R\,\uR\Bigr).\label{eq:R-mod}
\end{align}
\end{definicion}

% Figura: punto fuente, punto campo y vector que los une.
\begin{figure}[ht]
\centering
\tdplotsetmaincoords{70}{118}
\begin{tikzpicture}[tdplot_main_coords, scale=1.55]
  \def\sx{1.5}\def\sy{-0.6}\def\sz{2.0}   % punto fuente r'
  \def\fx{1.4}\def\fy{2.4}\def\fz{0.4}   % punto campo r
  \draw[eje] (0,0,0) -- (3.0,0,0) node[below left] {\small $x$};
  \draw[eje] (0,0,0) -- (0,3.1,0) node[right] {\small $y$};
  \draw[eje] (0,0,0) -- (0,0,3.1) node[above] {\small $z$};
  % vectores de posición
  \draw[->, line width=1.4pt, azulacero] (0,0,0) -- (\sx,\sy,\sz);
  \draw[->, line width=1.4pt, azulprof] (0,0,0) -- (\fx,\fy,\fz);
  % vector relativo R = r - r'
  \draw[->, line width=1.8pt, rojoteo] (\sx,\sy,\sz) -- (\fx,\fy,\fz);
  % versor u_R en el punto campo
  \pgfmathsetmacro{\Rx}{\fx-\sx}\pgfmathsetmacro{\Ry}{\fy-\sy}\pgfmathsetmacro{\Rz}{\fz-\sz}
  \pgfmathsetmacro{\Rn}{sqrt(\Rx*\Rx+\Ry*\Ry+\Rz*\Rz)}
  \draw[->, line width=1.6pt, verdeteo] (\fx,\fy,\fz) -- ++({0.95*\Rx/\Rn},{0.95*\Ry/\Rn},{0.95*\Rz/\Rn});
  \fill[azulacero] (\sx,\sy,\sz) circle (1.8pt);
  \fill[azulprof] (\fx,\fy,\fz) circle (1.8pt);
  \fill (0,0,0) circle (1.2pt);
  % etiquetas
  \node[azulacero, font=\normalsize] at (0.35,-0.55,1.0) {$\vc{r}\,'$};
  \node[azulprof, font=\normalsize] at (0.55,1.35,-0.15) {$\vc{r}$};
  \node[rojoteo, font=\normalsize, fill=white, fill opacity=0.8, text opacity=1, inner sep=1pt]
     at (1.55,1.15,1.5) {$\vc{R}=\vc{r}-\vc{r}\,'$};
  \node[verdeteo, font=\normalsize] at ({\fx+1.3*\Rx/\Rn},{\fy+1.3*\Ry/\Rn},{\fz+1.3*\Rz/\Rn}) {$\uR$};
  \node[font=\small, text=black!80, anchor=east] at (\sx+0.05,\sy-0.05,\sz+0.1) {fuente};
  \node[font=\small, text=black!80, anchor=north] at (\fx,\fy,\fz-0.12) {campo};
\end{tikzpicture}
\caption{Funciones de la posición relativa. El vector $\vc{r}\,'$ (azul claro)
localiza el \emph{punto fuente}, $\vc{r}$ (azul oscuro) el \emph{punto campo}, y
$\vc{R}=\vc{r}-\vc{r}\,'$ (rojo) apunta de la fuente al punto donde se evalúa
el campo; $\uR=\vc{R}/R$ (verde) es su versor.}
\label{fig:posicion-relativa}
\end{figure}


Interesan campos escalares $f(\vc{R})$ y vectoriales $\vc{A}(\vc{R})$ que dependen
de $\vc{r}$ y $\vc{r}\,'$ \emph{solo} a través de $\vc{R}$. Hay ahora dos
operadores nabla, que actúan sobre variables distintas:
\begin{equation}\label{eq:nabla-nabla-prima}
  \nab=\pd{}{x}\,\ux+\pd{}{y}\,\uy+\pd{}{z}\,\uz
  \qquad\text{y}\qquad
  \nab'=\pd{}{x'}\,\ux+\pd{}{y'}\,\uy+\pd{}{z'}\,\uz ,
\end{equation}
las derivadas respecto de las coordenadas del punto campo y del punto fuente,
respectivamente. Se diferencian en el signo, como se aprecia ya en la primera
derivada parcial:
\begin{equation}\label{eq:dR}
  \pd{R}{x}=\frac{1}{2}\,\frac{2(x-x')}{R}=\frac{R_x}{R},
  \qquad\text{mientras que}\qquad
  \pd{R}{x'}=-\frac{R_x}{R}.
\end{equation}

\begin{lema}[$\nab$ y $\nab'$ difieren en un signo]\label{lema:nabla-prima}
Si $g$ depende de $\vc{r}$ y $\vc{r}\,'$ solo a través de $\vc{R}=\vc{r}-\vc{r}\,'$,
entonces $\nab g=-\nab'g$ (y análogamente para la divergencia y el rotacional de
campos vectoriales de $\vc{R}$).
\end{lema}
\begin{proof}
Como $g$ es función de $x-x'$, la regla de la cadena da
$\partial g/\partial x'=-\partial g/\partial x$, y lo mismo para $y$ y $z$.
\end{proof}

% -------------------------------------------------------------------------
\subsection{Derivadas de potencias de $R$}\label{sec:derivadas-R}
% -------------------------------------------------------------------------
\begin{proposicion}[Gradiente de $R^{n}$]\label{prop:grad-Rn}
Para cualquier $n$ y $R\neq0$,
\begin{equation}\label{eq:grad-Rn}
  \nab(R^{n})=-\nab'(R^{n})=n\,R^{n-1}\,\uR=n\,R^{n-2}\,\vc{R}.
\end{equation}
\end{proposicion}
\begin{proof}
Por la regla de la cadena y la \cref{eq:dR}, cada componente es
\[
  \pd{R^{n}}{x}=\pd{R^{n}}{R}\,\pd{R}{x}=n\,R^{n-1}\,\frac{R_x}{R}=n\,R^{n-2}\,R_x ,
\]
y análogamente para $y,z$. Agrupando, $\nab(R^{n})=nR^{n-2}\bigl(R_x\,\ux+R_y\,\uy+R_z\,\uz\bigr)=nR^{n-2}\vc{R}=nR^{n-1}\uR$.
La igualdad con $-\nab'$ es el \cref{lema:nabla-prima}.
\end{proof}

\begin{proposicion}[Divergencia y rotacional de $R^{n}\uR$]\label{prop:div-rot-Rn}
Para $R\neq0$ y cualquier $n$ real (en particular entero),
\begin{equation}\label{eq:div-rot-Rn}
  \nab\cdot\bigl(R^{n}\,\uR\bigr)=-\nab'\cdot\bigl(R^{n}\,\uR\bigr)=(n+2)\,R^{n-1},
  \qquad
  \nab\times\bigl(R^{n}\,\uR\bigr)=-\nab'\times\bigl(R^{n}\,\uR\bigr)=\vc{0}.
\end{equation}
En particular, para $n=-2$: $\nab\cdot(\uR/R^{2})=0$ si $R\neq0$.
\end{proposicion}
\begin{proof}
Escribimos $R^{n}\uR=R^{n-1}\vc{R}$. Como $\nab\cdot\vc{R}=\partial_x(x-x')+\partial_y(y-y')+\partial_z(z-z')=3$ y
$\nab\times\vc{R}=\vc{0}$ (cada componente de $\vc{R}$ solo depende de su propia
coordenada), las reglas del producto \eqref{eq:prod-c} y \eqref{eq:prod-e} y la
\cref{prop:grad-Rn} dan
\begin{align*}
  \nab\cdot\bigl(R^{n-1}\vc{R}\bigr)&=R^{n-1}\,\nab\cdot\vc{R}+\vc{R}\cdot\nab R^{n-1}
    =3R^{n-1}+(n-1)R^{n-2}\,\vc{R}\cdot\uR=(n+2)\,R^{n-1},\\
  \nab\times\bigl(R^{n-1}\vc{R}\bigr)&=\nab R^{n-1}\times\vc{R}+R^{n-1}\,\nab\times\vc{R}
    =(n-1)R^{n-2}\,\uR\times\vc{R}+\vc{0}=\vc{0},
\end{align*}
pues $\uR\parallel\vc{R}$. Las igualdades con $\nab'$ son el \cref{lema:nabla-prima}.
En el \cref{ejer:div-Rn} se repite el cálculo de la divergencia componente a
componente.
\end{proof}

\begin{observacion}[Comprobaciones y el caso $n=-2$]\label{obs:casos-n}
Para $n=1$, $R\,\uR=\vc{R}$ y $\nab\cdot\vc{R}=3$, como en el \cref{ej:div}; para
$n=0$, $\nab\cdot\uR=2/R$; y para $n=-2$ el factor $(n+2)$ se anula. Pero
atención: el campo $\uR/R^{2}$ \emph{no está definido} en $R=0$ y su divergencia,
nula en todo punto $R\neq0$, tiene una contribución \emph{singular} concentrada
en el origen que la \cref{eq:div-rot-Rn} no recoge. Esa contribución es una delta de
Dirac (\cref{teo:delta}).
\end{observacion}

% -------------------------------------------------------------------------
\subsection{La delta de Dirac}\label{sec:delta}
% -------------------------------------------------------------------------
\begin{definicion}[Delta de Dirac]\label{def:delta}
La \emph{delta de Dirac} unidimensional $\delta(x)$ es el objeto que cumple
\begin{equation}\label{eq:delta1d}
  \delta(x)=\begin{cases}0,&x\neq0\\ \infty,&x=0\end{cases},
  \qquad
  \int_{-\infty}^{\infty}\delta(x)\,\dif x=1,
  \qquad
  \int_{-\infty}^{\infty}f(x)\,\delta(x-a)\,\dif x=f(a)
\end{equation}
para toda función continua $f$ (propiedad de selección o «muestreo»). En tres
dimensiones, $\delta^{3}(\vc{r})=\delta(x)\,\delta(y)\,\delta(z)$ cumple
\begin{equation}\label{eq:delta3d}
  \delta^{3}(\vc{r})=\begin{cases}0,&\vc{r}\neq\vc{0}\\ \infty,&\vc{r}=\vc{0}\end{cases},
  \qquad
  \int_{V_\infty}\delta^{3}(\vc{r})\,\dV=1,
  \qquad
  \int_{V_\infty}f(\vc{r})\,\delta^{3}(\vc{r}-\vc{a})\,\dV=f(\vc{a}),
\end{equation}
donde $V_\infty$ es todo el espacio (o cualquier región que contenga a
$\vc{a}$; si no lo contiene, la integral es nula).
\end{definicion}

\begin{observacion}[Qué es en realidad]\label{obs:delta}
La delta no es una función en el sentido usual (ninguna función es nula salvo
en un punto y tiene integral~$1$), sino una \emph{distribución}: un objeto que
solo tiene sentido dentro de una integral, donde actúa como un «detector» de
valores puntuales. Intuitivamente es el límite de picos cada vez más estrechos y
altos de área unidad, por ejemplo gaussianas
$\delta_\sigma(x)=\tfrac{1}{\sigma\sqrt{2\pi}}\,\eul^{-x^{2}/2\sigma^{2}}$ con
$\sigma\to0$ (\cref{fig:delta}). Representa matemáticamente una carga (o masa)
puntual: densidad nula en todas partes menos en un punto, y carga total finita.
\end{observacion}

% Figura: aproximaciones de la delta de Dirac por gaussianas de área unidad.
\begin{figure}[ht]
\centering
\begin{tikzpicture}
  \begin{axis}[tema, width=0.62\textwidth, height=5.6cm, xlabel={$x$}, ylabel={$\delta_\sigma(x)$},
      xmin=-2.2, xmax=2.2, ymin=0, ymax=3.6, samples=200, domain=-2.2:2.2,
      legend pos=north east, legend cell align=left, grid=major]
    \addplot[azulclaro!80!azulacero, line width=1.1pt] {1/(1*sqrt(2*pi))*exp(-x^2/(2*1^2))};
    \addplot[azulacero, line width=1.1pt] {1/(0.5*sqrt(2*pi))*exp(-x^2/(2*0.5^2))};
    \addplot[azulprof, line width=1.1pt] {1/(0.25*sqrt(2*pi))*exp(-x^2/(2*0.25^2))};
    \addplot[rojoteo, line width=1.1pt, samples=400] {1/(0.12*sqrt(2*pi))*exp(-x^2/(2*0.12^2))};
    \legend{$\sigma=1$, $\sigma=0{,}5$, $\sigma=0{,}25$, $\sigma=0{,}12$}
  \end{axis}
\end{tikzpicture}
\caption{La delta de Dirac como límite de picos cada vez más estrechos y altos
con área unidad: $\delta_\sigma(x)=\frac{1}{\sigma\sqrt{2\pi}}\,\eul^{-x^{2}/2\sigma^{2}}$.
Cuando $\sigma\to0$, $\int\delta_\sigma\,\dif x=1$ se conserva y el pico se
concentra en $x=0$.}
\label{fig:delta}
\end{figure}


\begin{teorema}[Divergencia del campo coulombiano y laplaciano de $1/R$]\label{teo:delta}
En el sentido de distribuciones,
\begin{equation}\label{eq:delta-div}
  \nab\cdot\Bigl(\frac{\uR}{R^{2}}\Bigr)=-\nab'\cdot\Bigl(\frac{\uR}{R^{2}}\Bigr)=4\pi\,\delta^{3}(\vc{R}),
  \qquad
  \Lap\Bigl(\frac{1}{R}\Bigr)=\Lapp\Bigl(\frac{1}{R}\Bigr)=-4\pi\,\delta^{3}(\vc{R}).
\end{equation}
\end{teorema}
\begin{proof}
\emph{Primera igualdad.} Por la \cref{prop:div-rot-Rn}, $\nab\cdot(\uR/R^{2})=0$ para
$R\neq0$: la distribución, si no es nula, está concentrada en $R=0$. Calculemos
su integral sobre una región $V$ que contenga a la fuente. Quitamos una bola
pequeña $B_\varepsilon$ de radio $\varepsilon$ alrededor de $\vc{r}\,'$; en
$V\setminus B_\varepsilon$ el campo es de clase $C^{1}$ y de divergencia nula, así
que por el teorema de Gauss (\cref{teo:gauss})
\[
  0=\iiint_{V\setminus B_\varepsilon}\nab\cdot\frac{\uR}{R^{2}}\,\dV
   =\oiint_{\partial V}\frac{\uR}{R^{2}}\cdot\dS-\oiint_{\partial B_\varepsilon}\frac{\uR}{R^{2}}\cdot\dS ,
\]
donde la normal exterior de la bola es $\uR$ (con el signo $-$ porque, para
$V\setminus B_\varepsilon$, la normal es la interior a la bola). Sobre
$\partial B_\varepsilon$, $\uR\cdot\uR/\varepsilon^{2}=1/\varepsilon^{2}$ es constante y
el flujo vale $\tfrac{1}{\varepsilon^{2}}\cdot4\pi\varepsilon^{2}=4\pi$. Por tanto
el flujo del campo por \emph{cualquier} superficie cerrada que rodee la fuente
es $4\pi$, esto es, la integral de su divergencia sobre cualquier región que
contenga a $\vc{r}\,'$ vale $4\pi$, independientemente del tamaño y la forma: es
exactamente la propiedad $\int\delta^{3}\,\dV=1$ multiplicada por $4\pi$.
Si $\varphi$ es una función continua cualquiera, el mismo argumento con
$\varepsilon\to0$ da $\int\varphi\,\nab\cdot(\uR/R^{2})\,\dV=4\pi\varphi(\vc{r}\,')$, que es
\eqref{eq:delta3d}.

\emph{Segunda igualdad.} Por la \cref{prop:grad-Rn} con $n=-1$,
$\nab(1/R)=-\uR/R^{2}$, de modo que $\Lap(1/R)=\nab\cdot\nab(1/R)=-\nab\cdot(\uR/R^{2})=-4\pi\,\delta^{3}(\vc{R})$.
Finalmente $\Lapp(1/R)=\nab'\cdot\nab'(1/R)=(-\nab)\cdot(-\nab)(1/R)=\Lap(1/R)$.
\end{proof}

En electrostática esta identidad es la ley de Gauss en forma local: el campo de
una carga puntual es proporcional a $\uR/R^{2}$ y su divergencia es una delta
centrada en la carga.

% -------------------------------------------------------------------------
\subsection{Polinomios de Legendre y desarrollo de $1/|\vc{r}-\vc{r}\,'|$}\label{sec:legendre}
% -------------------------------------------------------------------------
El desarrollo del inverso de la distancia en serie de potencias de
$r'/r$ es la base de la expansión multipolar. Sus coeficientes son los
polinomios de Legendre.

\begin{definicion}[Polinomios de Legendre]\label{def:legendre}
Los polinomios de Legendre $P_n(x)$, $x\in[-1,1]$, vienen dados por la fórmula de
Rodrigues $P_n(x)=\frac{1}{2^{n}n!}\frac{\dif^{n}}{\dif x^{n}}(x^{2}-1)^{n}$. Los
primeros son
\[
  P_0(x)=1,\qquad P_1(x)=x,\qquad P_2(x)=\tfrac{1}{2}(3x^{2}-1),\qquad P_3(x)=\tfrac{1}{2}(5x^{3}-3x),
\]
y cumplen $P_n(1)=1$ y la recurrencia
$(n+1)P_{n+1}=(2n+1)\,x\,P_n-n\,P_{n-1}$.
\end{definicion}

\begin{teorema}[Desarrollo en polinomios de Legendre]\label{teo:legendre}
Si $r'<r$ y $\theta'$ es el ángulo que forman $\vc{r}$ y $\vc{r}\,'$,
\begin{equation}\label{eq:legendre}
  \frac{1}{\abs{\vc{r}-\vc{r}\,'}}=\frac{1}{r}\sum_{n=0}^{\infty}\Bigl(\frac{r'}{r}\Bigr)^{n}P_n(\cos\theta').
\end{equation}
\end{teorema}
\begin{proof}
Por el teorema del coseno,
$\abs{\vc{r}-\vc{r}\,'}^{2}=r^{2}+r'^{2}-2rr'\cos\theta'=r^{2}\bigl(1-2xt+t^{2}\bigr)$
con $t=r'/r<1$ y $x=\cos\theta'$. Por tanto
$\frac{1}{\abs{\vc{r}-\vc{r}\,'}}=\frac{1}{r}(1+u)^{-1/2}$ con $u=t^{2}-2xt$ pequeño,
y el desarrollo binomial $(1+u)^{-1/2}=1-\tfrac{u}{2}+\tfrac{3u^{2}}{8}-\tfrac{5u^{3}}{16}+\cdots$
ordenado por potencias de $t$ da
\[
  1+x\,t+\tfrac{1}{2}(3x^{2}-1)\,t^{2}+\tfrac{1}{2}(5x^{3}-3x)\,t^{3}+\cdots
  =P_0+P_1\,t+P_2\,t^{2}+P_3\,t^{3}+\cdots ,
\]
es decir, la función generatriz $(1-2xt+t^{2})^{-1/2}=\sum_nP_n(x)\,t^{n}$ de los
polinomios de Legendre~\cite{ARFKEN}.
\end{proof}

% Figura: polinomios de Legendre P0..P3.
\begin{figure}[ht]
\centering
\begin{tikzpicture}
  \begin{axis}[tema, width=0.62\textwidth, height=6.2cm, xlabel={$x=\cos\theta'$}, ylabel={$P_n(x)$},
      xmin=-1, xmax=1, ymin=-1.05, ymax=1.15, samples=100, domain=-1:1,
      xtick={-1,-0.5,0,0.5,1}, legend pos=south east, legend cell align=left]
    \addplot[grispizarra, line width=1.1pt] {1};
    \addplot[azulclaro!70!azulacero, line width=1.1pt] {x};
    \addplot[azulacero, line width=1.1pt] {0.5*(3*x^2-1)};
    \addplot[azulprof, line width=1.1pt] {0.5*(5*x^3-3*x)};
    \legend{$P_0$, $P_1$, $P_2$, $P_3$}
  \end{axis}
\end{tikzpicture}
\caption{Primeros polinomios de Legendre en $[-1,1]$. Todos valen $1$ en $x=1$
(dirección de $\vc{r}\,'$ paralela a $\vc{r}$).}
\label{fig:legendre}
\end{figure}


Cada término tiene un significado físico: $n=0$ es el campo de un monopolo
($\propto1/r$); $n=1$, un dipolo ($\propto r'\cos\theta'/r^{2}$), etc. Si $r'>r$
basta intercambiar los papeles de $r$ y $r'$.
